% ============================================================
% Figure 4 caption -- local mirror, paste into Overleaf manuscript.tex
% ============================================================
% Figure 4 is assembled from two separately-rendered files (LaTeX/Overleaf
% side, not composited in R):
%   figs/revision_2026/Figure4_network_trio.pdf   -- panels (A)-(C), now
%     8.8x2.9in (tightened from 12x4.3in in the 2026-08-25 design pass),
%     with orange phantom-edge highlighting in panels B/C.
%   figs/revision_2026/Figure4_metric_strip.pdf    -- panels (D)-(E)
% Suggested layout: network trio as the top row (full width), metric
% strip as the bottom row (centered or full width), inside one figure
% environment with subfloats, or two \includegraphics calls stacked
% vertically under one \caption.
% ============================================================

\caption{\textbf{Omitting slow context makes estimated symptom networks look more connected.}
(A) True symptom--symptom coupling matrix used to generate the data.
(B) Network estimated from pooled symptom profiles when the slow contextual field varies across people but is omitted from the estimator. Orange edges mark phantom couplings: connections absent in the true network that nonetheless appear in this estimate.
(C) Network estimated from the same pooled data after including the slow contextual field as a covariate; the same phantom-edge coloring is applied, and few or no orange edges remain.
(D) Omitting \(P\) produces the largest excess estimated global strength relative to the true coupling matrix.
(E) Error on truly absent edges is also largest when \(P\) is omitted, indicating more apparent coupling among symptom pairs with no true direct edge. Points and intervals in panels D--E show means and standard errors across 30 independent replicates.}
